\section{Objectives}

SmartExam has one overarching goal: replace the fragmented, manual workflow of practical lab exams with a single, auditable digital platform. Below are the six concrete objectives that define how we achieve this.

\textbf{Objective 1 — Centralise administration.} Scheduling, user management, lab allocation, and session configuration shall all happen within one platform. There should be no separate spreadsheets, no external booking tools, and no email chains to negotiate room availability.

\textbf{Objective 2 — Strengthen live visibility.} Invigilators shall see every active session on one dashboard. Real-time focus timelines, active process lists, and configurable alert thresholds replace the current model of physically walking between rows of students \cite{rbac1}.

\textbf{Objective 3 — Build a trustworthy evidence chain.} Every submission shall carry a timestamp, a hardware identity (HWID), and a linked event log from that session. If a student disputes a result, instructors can reconstruct exactly what happened at that workstation during the exam.

\textbf{Objective 4 — Enforce procedural fairness.} Access is role-gated and hardware-bound. Students can only log in from their registered machine. The exam client runs in fullscreen lockdown mode. No student gets an advantage from running unauthorized software.

\textbf{Objective 5 — Reduce administrative overhead.} Recurring manual tasks — attendance recording, file collection, result tabulation — shall be automated wherever possible. Teachers get a consolidated post-exam report rather than scattered data sources.

\textbf{Objective 6 — Deliver within FYP constraints.} The system shall be technically feasible within a two-semester FYP timeline, academically defensible, and structured so that future cohorts can extend individual modules without rebuilding the core.
